\documentclass[12pt]{article}
\usepackage{graphicx} 
\usepackage{amsmath, amssymb,amsthm}
\newtheorem{remark}{Remark}
\usepackage[margin=1in]{geometry}
\usepackage{hyperref}

\title{ECCV 2026}
\author{Notes for Abhinay}
\date{February 2026}

\newtheorem{lemma}{Lemma}
\newtheorem{proposition}{Proposition}
\documentclass[12pt]{article}
\usepackage{graphicx}
\usepackage{amsmath, amssymb, amsthm}
\usepackage[margin=1in]{geometry}
\usepackage{hyperref}
\usepackage{booktabs}
\usepackage{algorithm}
\usepackage{algorithmic}

\title{Structure-Aware Generation via Representation Diffusion:\\
Theoretical Foundations}
\author{Abhinay Belde}
\date{February 2026}

\newtheorem{theorem}{Theorem}[section]


\newtheorem{corollary}[theorem]{Corollary}
\newtheorem{definition}[theorem]{Definition}

\newtheorem{assumption}[theorem]{Assumption}

\DeclareMathOperator{\Tr}{Tr}
\DeclareMathOperator{\diag}{diag}
\DeclareMathOperator{\softmax}{softmax}

\begin{document}

\maketitle

\tableofcontents

\section{Theoretical Viewpoint}

\subsection{Two-Stage Modeling in Representation Space}

Let \(X\) be a random variable of natural images with distribution \(p_X\).
Let \(E_{\mathrm{g}} : \mathcal{X} \to \mathbb{R}^{d_g}\) be the I-JEPA encoder and
\(E_{\mathrm{s}} : \mathcal{X} \to \mathbb{R}^{N \times d_s}\) be the SAM-based encoder
that outputs per-mask embeddings. For each image \(x\), we define
\[
G = E_{\mathrm{g}}(X) \in \mathbb{R}^{d_g}, \qquad
S = E_{\mathrm{s}}(X) \in \mathbb{R}^{N \times d_s}.
\]
This induces a joint representation distribution
\[
p_{G,S}(g,s)
= \int p_X(x)\,\delta\big(g - E_{\mathrm{g}}(x)\big)\,\delta\big(s - E_{\mathrm{s}}(x)\big)\,dx,
\]
where \(\delta\) is the Dirac delta.

Our generative goal can be written as
\[
p_X(x)
= \int p_{X \mid G,S}(x \mid g,s)\,p_{G,S}(g,s)\,dg\,ds,
\]
where \(p_{G,S}\) is the distribution of self-supervised representations and
\(p_{X \mid G,S}\) is a decoder from representation space back to pixels.
Our two-stage model approximates these components as
\(\hat{p}_{G,S}\) (Stage~1) and \(\hat{p}_{X \mid G,S}\) (Stage~2),
yielding the overall generative model
\[
\hat{p}_X(x)
= \int \hat{p}_{X \mid G,S}(x \mid g,s)\,\hat{p}_{G,S}(g,s)\,dg\,ds.
\]

\subsection{Stage 1 as Rectified Flow / Diffusion}

Let \(Z = (G,S)\) be the concatenated representation, taking values in
\(\mathbb{R}^d\) with \(d = d_g + N d_s\).
We choose a simple base distribution \(p_0 = \mathcal{N}(0, I_d)\)
and denote the data distribution in representation space by
\(p_1 = p_Z = p_{G,S}\).
In the rectified-flow viewpoint, we consider a family of intermediate
distributions \(\{p_t\}_{t \in [0,1]}\) and a time-dependent
velocity field \(v_\theta : \mathbb{R}^d \times [0,1] \to \mathbb{R}^d\)
such that the flow \(\phi_t\) satisfies
\[
\frac{d}{dt}\phi_t(z_0) = v_\theta\big(\phi_t(z_0), t\big), \qquad
\phi_0(z_0) = z_0, \quad z_0 \sim p_0,
\]
and \(\phi_{1\sharp} p_0 \approx p_1\), where \(\phi_{1\sharp}\) denotes the
pushforward of measures.

In practice, we parameterize this flow via a diffusion-style noise-adding
process and a denoising network.
Let \(Z \sim p_1\) and \(\varepsilon \sim \mathcal{N}(0, I_d)\).
For a noise schedule \(\{\alpha_t\}_{t \in [0,1]}\), we define
\[
Z_t = \sqrt{\alpha_t}\, Z + \sqrt{1 - \alpha_t}\,\varepsilon.
\]
We train a network \(f_\theta(Z_t, t)\) to predict the clean
representations \(Z\) from \(Z_t\) via the objective
\[
\mathcal{L}_{\mathrm{denoise}}(\theta)
= \mathbb{E}_{Z, t, \varepsilon}
\big[\|f_\theta(Z_t, t) - Z\|_2^2\big].
\]
Under standard assumptions on the noise schedule,
this denoising objective corresponds to learning the score or velocity
field associated with a flow that transports \(p_0\) to \(p_1\).

In our setting, the network has structured outputs
\[
f_\theta(Z_t, t)
= \big(f_\theta^{\mathrm{g}}(Z_t, t), f_\theta^{\mathrm{s}}(Z_t, t)\big),
\]
corresponding to predictions of the global embedding \(G\) and the
segmentation embeddings \(S\).
The loss decomposes as
\[
\mathcal{L}_{\mathrm{global}}
= \mathbb{E}\big[\|f_\theta^{\mathrm{g}}(Z_t, t) - G\|_2^2\big], \qquad
\mathcal{L}_{\mathrm{seg}}
= \mathbb{E}\big[\|f_\theta^{\mathrm{s}}(Z_t, t) - S\|_2^2\big],
\]
so that the learned flow respects the joint structure of
\((G,S)\) in representation space.

\subsection{Stage 2 as Conditional Decoder}

Stage~2 is a conditional decoder from representation space to pixels.
We define a generator
\[
D_\psi : \mathbb{R}^{d_g} \times \mathbb{R}^{N \times d_s} \to \mathcal{X}
\]
that induces a conditional distribution
\[
\hat{p}_{X \mid G,S}(x \mid g,s) = p_\psi(x \mid g,s).
\]
We train this decoder with a conditional GAN objective and auxiliary
reconstruction losses.
Let \(D_\phi\) be a discriminator.
We optimize
\begin{align*}
\min_\psi \max_\phi \quad
&\mathbb{E}_{(x,g,s) \sim (X,G,S)}
\big[\log D_\phi(x,g,s)\big] \\
&+ \mathbb{E}_{(g,s) \sim p_{G,S},
\ \tilde{x} \sim p_\psi(\cdot \mid g,s)}
\big[\log(1 - D_\phi(\tilde{x}, g,s))\big]
+ \lambda_{\mathrm{rec}}\,\mathcal{L}_{\mathrm{rec}}(\psi),
\end{align*}
where \(\mathcal{L}_{\mathrm{rec}}\) combines perceptual and pixel-level losses
between \(\tilde{x} = D_\psi(g,s)\) and the observed image \(x\).

At inference, the overall sampling procedure is:
\begin{enumerate}
    \item Sample \(z_0 \sim \mathcal{N}(0, I_d)\).
    \item Evolve \(z_0\) through the learned flow (or reverse diffusion)
          to obtain \(\hat{Z}_1 = (\hat{G}, \hat{S})\).
    \item Decode an image \(\hat{X} = D_\psi(\hat{G}, \hat{S})\).
\end{enumerate}
This realizes sampling from
\[
\hat{X} \sim \int \hat{p}_{X \mid G,S}(x \mid g,s)\,
\hat{p}_{G,S}(g,s)\,dg\,ds.
\]

\subsection{Alignment and Diversity Losses Prevent Representation Collapse}

For a fixed image \(x\), let its per-mask embeddings be
\[
S(x) = \big(s_1(x), \dots, s_N(x)\big), \qquad
s_i(x) \in \mathbb{R}^{d_s}.
\]
Define the sample mean
\[
\bar{s}(x) = \frac{1}{N} \sum_{i=1}^N s_i(x),
\]
and the sample covariance
\[
C(x) = \frac{1}{N} \sum_{i=1}^N
\big(s_i(x) - \bar{s}(x)\big)\big(s_i(x) - \bar{s}(x)\big)^\top
\in \mathbb{R}^{d_s \times d_s}.
\]

\paragraph{Diversity loss.}
We introduce a diversity-promoting term
\[
\mathcal{L}_{\mathrm{div}}(x)
= - \log \det \big(C(x) + \epsilon I\big),
\]
with a small \(\epsilon > 0\) for numerical stability.
If all mask embeddings collapse to a single vector, i.e.,
\(s_i(x) = s_j(x)\) for all \(i,j\), then \(C(x) = 0\) and
\(\det\big(C(x) + \epsilon I\big) = \epsilon^{d_s}\).
In the limit \(\epsilon \to 0\),
the determinant is zero and \(\mathcal{L}_{\mathrm{div}}(x) \to +\infty\).

More generally, let \(\lambda_1, \dots, \lambda_{d_s}\) be the eigenvalues
of \(C(x)\).
Then
\[
\det\big(C(x) + \epsilon I\big)
= \prod_{k=1}^{d_s} (\lambda_k + \epsilon),
\qquad
\mathcal{L}_{\mathrm{div}}(x)
= - \sum_{k=1}^{d_s} \log (\lambda_k + \epsilon).
\]
Any eigenvalue \(\lambda_k \to 0\) causes the corresponding term
to blow up.
Minimizing \(\mathcal{L}_{\mathrm{div}}\) therefore encourages the
covariance \(C(x)\) to be full-rank and well-conditioned, i.e.,
it discourages low-dimensional collapse of mask embeddings.

\paragraph{Alignment loss.}
Let \(g(x)\) denote the global embedding of image \(x\).
For a batch \(\{x^{(b)}\}_{b=1}^B\), a contrastive alignment loss can be
written as
\[
\mathcal{L}_{\mathrm{align}} =
\mathbb{E}\bigg[
- \log
\frac{\exp\big(\mathrm{sim}(g(x), s_i(x)) / \tau\big)}
{\sum_{b',j} \exp\big(\mathrm{sim}(g(x), s_j(x^{(b')})) / \tau\big)}
\bigg],
\]
where \(\mathrm{sim}(\cdot,\cdot)\) is a similarity function
(e.g., dot product or cosine) and \(\tau > 0\) is a temperature.

If all mask embeddings collapse to a constant vector \(c\) across images
and masks, i.e., \(s_i(x) = c\) for all \(i,x\), and the global embeddings
also collapse or trivially align with \(c\), then
\(\mathrm{sim}(g(x), s_i(x))\) becomes nearly independent of
\((x,i)\).
In this regime, the numerator and denominator of the contrastive ratio
are similar across all pairs, so the loss cannot effectively
distinguish positive from negative pairs and remains high.
Minimizing \(\mathcal{L}_{\mathrm{align}}\) thus encourages
\((g(x), s_i(x))\) to retain image-specific information and to be
mutually informative, preventing trivial collapse.

\paragraph{Joint effect.}
The total embedding-related loss combines reconstruction, alignment,
and diversity terms, e.g.,
\[
\mathcal{L}_{\mathrm{emb}}
= \lambda_2 \mathcal{L}_{\mathrm{seg}}
+ \lambda_3 \mathcal{L}_{\mathrm{align}}
+ \lambda_4 \mathcal{L}_{\mathrm{div}}.
\]
Any configuration where the mask embeddings collapse to a single
structural code leads to:
(i) a large diversity loss due to nearly singular covariance,
(ii) a large alignment loss due to poor discrimination between positive
and negative pairs, and
(iii) increased reconstruction error relative to the teacher-forced SAM
embeddings.
Therefore, minimizing \(\mathcal{L}_{\mathrm{emb}}\) drives the model
towards non-degenerate, diverse mask embeddings that remain coherent
with global representations while avoiding trivial collapse.
%\subsection{Alignment and Diversity Losses Prevent Representation Collapse}

\begin{lemma}[Diversity loss penalizes collapsed embeddings]
Let \(S(x) = (s_1(x),\dots,s_N(x))\) with \(s_i(x) \in \mathbb{R}^{d_s}\), and let
\[
\bar{s}(x) = \frac{1}{N} \sum_{i=1}^N s_i(x), \qquad
C(x) = \frac{1}{N} \sum_{i=1}^N (s_i(x) - \bar{s}(x))(s_i(x) - \bar{s}(x))^\top.
\]
Assume the diversity loss
\[
\mathcal{L}_{\mathrm{div}}(x)
= - \log \det \big(C(x) + \epsilon I\big), \quad \epsilon > 0.
\]
If all mask embeddings collapse to a single vector, i.e.\ \(s_i(x) = c\) for all \(i\),
then \(C(x) = 0\) and therefore
\[
\mathcal{L}_{\mathrm{div}}(x) = - d_s \log \epsilon.
\]
Moreover, for any non-collapsed configuration with eigenvalues
\(\lambda_1,\dots,\lambda_{d_s}\) of \(C(x)\),
\[
\mathcal{L}_{\mathrm{div}}(x)
= - \sum_{k=1}^{d_s} \log (\lambda_k + \epsilon)
\]
strictly decreases as any \(\lambda_k\) increases from \(0\).
In particular, any configuration in which one or more eigenvalues
\(\lambda_k\) are close to zero yields a large value of
\(\mathcal{L}_{\mathrm{div}}(x)\), so the loss strongly penalizes
collapsed or nearly collapsed spectra.
\end{lemma}

\begin{proof}
If \(s_i(x) = c\) for all \(i\), then \(\bar{s}(x) = c\) and
\(s_i(x) - \bar{s}(x) = 0\) for all \(i\), so \(C(x) = 0\).
Thus
\[
\det(C(x) + \epsilon I) = \det(\epsilon I) = \epsilon^{d_s},
\]
and therefore \(\mathcal{L}_{\mathrm{div}}(x) = - d_s \log \epsilon\).

In the general case, let \(\lambda_1,\dots,\lambda_{d_s}\) be the eigenvalues of
\(C(x)\), so
\[
\det(C(x) + \epsilon I) = \prod_{k=1}^{d_s} (\lambda_k + \epsilon),
\quad
\mathcal{L}_{\mathrm{div}}(x)
= - \sum_{k=1}^{d_s} \log (\lambda_k + \epsilon).
\]
For each \(k\),
\[
\frac{\partial \mathcal{L}_{\mathrm{div}}}{\partial \lambda_k}
= - \frac{1}{\lambda_k + \epsilon} < 0,
\]
so increasing any \(\lambda_k\) strictly decreases
\(\mathcal{L}_{\mathrm{div}}(x)\).
In particular, when some \(\lambda_k\) are close to zero, the terms
\(-\log(\lambda_k + \epsilon)\) are large, making
\(\mathcal{L}_{\mathrm{div}}(x)\) large.
Hence the diversity loss penalizes collapsed or nearly collapsed spectra.
\end{proof}


\begin{lemma}[No trivial joint minimizer under collapse]
Fix an image $x$ and let $S^\star(x) = (s^\star_1(x),\dots,s^\star_N(x))$
denote the corresponding teacher per-mask embeddings (e.g., from a
segmentation encoder such as SAM), with sample covariance
$C^\star(x)$ defined by
\[
\bar{s}^\star(x) = \frac{1}{N} \sum_{i=1}^N s^\star_i(x), \qquad
C^\star(x) = \frac{1}{N} \sum_{i=1}^N
\big(s^\star_i(x) - \bar{s}^\star(x)\big)\big(s^\star_i(x) - \bar{s}^\star(x)\big)^\top.
\]


Assume:
\begin{enumerate}
    \item[(i)] $C^\star(x)$ is non-zero for some $x$, i.e.\ the teacher SAM
    embeddings are not collapsed for at least one image;
    \item[(ii)] the segmentation reconstruction loss $\mathcal{L}_{\mathrm{seg}}$
    is minimized (over the model parameters) only when, for every image $x$,
    the predicted mask embeddings $S(x)$ equal the teacher embeddings
    $S^\star(x)$;
    \item[(iii)] the alignment loss $\mathcal{L}_{\mathrm{align}}$ is minimized
    only when the joint configuration of global and mask embeddings
    $\{(G(x), S(x))\}_x$ is discriminative across images (in particular,
    assigning the same constant mask embedding to all images does not minimize
    $\mathcal{L}_{\mathrm{align}}$);
    \item[(iv)] the diversity loss for an image $x$ is
    \[
    \mathcal{L}_{\mathrm{div}}(x)
    = - \log \det\big(C(x) + \epsilon I\big), \quad \epsilon > 0,
    \]
    where $C(x)$ is the empirical covariance of the predicted mask embeddings
    $S(x)$ defined analogously to $C^\star(x)$.
\end{enumerate}

Let
\[
\mathcal{L}_{\mathrm{emb}}
= \mathcal{L}_{\mathrm{seg}}
+ \lambda_3 \mathcal{L}_{\mathrm{align}}
+ \lambda_4 \mathcal{L}_{\mathrm{div}},
\qquad \lambda_3,\lambda_4 > 0.
\]
Then any configuration in which, for some image $x$, all predicted mask
embeddings collapse to a single constant vector, i.e.\ $s_i(x) = c$ for all $i$,
cannot be a global minimizer of $\mathcal{L}_{\mathrm{emb}}$.
\end{lemma}

\begin{proof}
Consider any configuration of model parameters for which there exists at least
one image $x_0$ such that all predicted mask embeddings collapse, i.e.\
\[
s_i(x_0) = c \quad \text{for all } i = 1,\dots,N
\]
for some $c \in \mathbb{R}^{d_s}$.

\medskip
\noindent\textbf{Step 1: Diversity loss is suboptimal under collapse.}
For this image $x_0$, the predicted embeddings satisfy
\[
\bar{s}(x_0) = \frac{1}{N} \sum_{i=1}^N s_i(x_0) = c,
\]
and hence $s_i(x_0) - \bar{s}(x_0) = 0$ for all $i$.
Thus the predicted covariance is
\[
C(x_0)
= \frac{1}{N} \sum_{i=1}^N
\big(s_i(x_0) - \bar{s}(x_0)\big)\big(s_i(x_0) - \bar{s}(x_0)\big)^\top
= 0.
\]
Therefore
\[
\mathcal{L}_{\mathrm{div}}(x_0)
= - \log \det\big(C(x_0) + \epsilon I\big)
= - \log \det(\epsilon I)
= - d_s \log \epsilon.
\]

Now, by assumption (i), for the teacher embeddings $S^\star(x_0)$ we have
$C^\star(x_0) \neq 0$, so at least one eigenvalue of $C^\star(x_0)$ is
strictly positive.
By the lemma on the diversity loss (proved previously), any configuration
with $C(x_0)$ closer to $C^\star(x_0)$ has strictly larger eigenvalues than
the collapsed case and thus a strictly smaller value of
$\mathcal{L}_{\mathrm{div}}(x_0)$.
In particular, setting $S(x_0) = S^\star(x_0)$ yields $C(x_0) = C^\star(x_0)$
and strictly decreases $\mathcal{L}_{\mathrm{div}}(x_0)$ compared to the
collapsed configuration.

\medskip
\noindent\textbf{Step 2: Segmentation reconstruction is suboptimal under collapse.}
By assumption (ii), the segmentation loss $\mathcal{L}_{\mathrm{seg}}$ is
minimized only when $S(x) = S^\star(x)$ for every image $x$.
For the image $x_0$ under the collapsed configuration $S(x_0) \equiv c$,
we have $S(x_0) \neq S^\star(x_0)$ (since $C^\star(x_0) \neq 0$ and hence
the teacher embeddings cannot all be equal).
Therefore, replacing $S(x_0)$ by $S^\star(x_0)$ (and keeping all other
parameters fixed) strictly reduces $\mathcal{L}_{\mathrm{seg}}$.

\medskip
\noindent\textbf{Step 3: Alignment loss is suboptimal under collapse.}
By assumption (iii), the alignment loss $\mathcal{L}_{\mathrm{align}}$ is
minimized only when the collection $\{(G(x), S(x))\}_x$ is discriminative
across images, and in particular a configuration where mask embeddings are
constant (or nearly constant) across images does not minimize
$\mathcal{L}_{\mathrm{align}}$.
Thus there exists at least one perturbation of the collapsed configuration
that decreases $\mathcal{L}_{\mathrm{align}}$.

\medskip
\noindent\textbf{Step 4: Strict improvement of the joint objective.}
Consider now moving from the collapsed configuration to a configuration where,
for the image $x_0$, we set $S(x_0) = S^\star(x_0)$ and adjust the model
parameters (if necessary) to reduce $\mathcal{L}_{\mathrm{align}}$ (as allowed
by (iii)), while leaving other images unchanged.
From Steps 1–3, this move strictly decreases at least one of
$\mathcal{L}_{\mathrm{seg}}$, $\mathcal{L}_{\mathrm{align}}$, or
$\mathcal{L}_{\mathrm{div}}$, and cannot increase the others arbitrarily
without undoing the improvements.
Since $\lambda_3,\lambda_4 > 0$, any strict decrease in one of these terms
(with the others not decreased negatively enough to compensate) produces a
strict decrease in
\[
\mathcal{L}_{\mathrm{emb}}
= \mathcal{L}_{\mathrm{seg}}
+ \lambda_3 \mathcal{L}_{\mathrm{align}}
+ \lambda_4 \mathcal{L}_{\mathrm{div}}.
\]

Hence the collapsed configuration cannot be a global minimizer of
$\mathcal{L}_{\mathrm{emb}}$.
\end{proof}

\begin{lemma}[SEG-RDM avoids collapsed mask embeddings]
Fix an image $x$ and let $S^\star(x) = (s^\star_1(x),\dots,s^\star_N(x))$
denote the corresponding teacher per-mask embeddings (e.g., from a
segmentation encoder such as SAM), with sample covariance
\[
\bar{s}^\star(x) = \frac{1}{N} \sum_{i=1}^N s^\star_i(x), \qquad
C^\star(x) = \frac{1}{N} \sum_{i=1}^N
\big(s^\star_i(x) - \bar{s}^\star(x)\big)\big(s^\star_i(x) - \bar{s}^\star(x)\big)^\top.
\]

Assume:
\begin{enumerate}
    \item[(i)] $C^\star(x)$ is non-zero for at least one image $x$, i.e.,
    the teacher embeddings are not collapsed for that image;
    \item[(ii)] the segmentation reconstruction loss $\mathcal{L}_{\mathrm{seg}}$
    is minimized (over model parameters) only when, for every image $x$,
    the predicted mask embeddings $S(x)$ equal the teacher embeddings
    $S^\star(x)$;
    \item[(iii)] the alignment loss $\mathcal{L}_{\mathrm{align}}$ is minimized
    only when the joint configuration $\{(G(x), S(x))\}_x$ is discriminative
    across images (in particular, assigning the same constant mask embedding
    to all images does not minimize $\mathcal{L}_{\mathrm{align}}$);
    \item[(iv)] for each image $x$, the diversity loss is
    \[
    \mathcal{L}_{\mathrm{div}}(x)
    = - \log \det\big(C(x) + \epsilon I\big), \quad \epsilon > 0,
    \]
    where $C(x)$ is the empirical covariance of the predicted mask embeddings
    $S(x)$ defined analogously to $C^\star(x)$.
\end{enumerate}

Let
\[
\mathcal{L}_{\mathrm{emb}}
= \mathcal{L}_{\mathrm{seg}}
+ \lambda_3 \mathcal{L}_{\mathrm{align}}
+ \lambda_4 \mathcal{L}_{\mathrm{div}},
\qquad \lambda_3,\lambda_4 > 0.
\]
Then any configuration in which, for some image $x$, all predicted mask
embeddings collapse to a single constant vector, i.e.\ $s_i(x) = c$ for all $i$,
cannot be a global minimizer of $\mathcal{L}_{\mathrm{emb}}$.
\end{lemma}

\begin{proof}
Suppose, for the sake of contradiction, that there exists a global minimizer
of $\mathcal{L}_{\mathrm{emb}}$ in which the predicted mask embeddings
collapse for at least one image $x_0$, i.e.\ there exists $c \in \mathbb{R}^{d_s}$
such that
\[
s_i(x_0) = c \quad \text{for all } i=1,\dots,N.
\]

\medskip
\noindent\textbf{Step 1: Diversity loss under collapse.}
For this image $x_0$, the predicted mean and covariance are
\[
\bar{s}(x_0) = \frac{1}{N} \sum_{i=1}^N s_i(x_0) = c,
\]
and therefore $s_i(x_0) - \bar{s}(x_0) = 0$ for all $i$, so
\[
C(x_0)
= \frac{1}{N} \sum_{i=1}^N
\big(s_i(x_0) - \bar{s}(x_0)\big)\big(s_i(x_0) - \bar{s}(x_0)\big)^\top
= 0.
\]
Hence
\[
\mathcal{L}_{\mathrm{div}}(x_0)
= - \log \det\big(C(x_0) + \epsilon I\big)
= - \log \det(\epsilon I)
= - d_s \log \epsilon.
\]

By assumption (i), for the teacher embeddings $S^\star(x_0)$ we have
$C^\star(x_0) \neq 0$, so at least one eigenvalue of $C^\star(x_0)$
is strictly positive.
As shown in the diversity lemma, any configuration with covariance
eigenvalues strictly larger than those of the collapsed case (which has
all eigenvalues equal to $0$) yields a strictly smaller value of
$\mathcal{L}_{\mathrm{div}}(x_0)$.
In particular, setting $S(x_0) = S^\star(x_0)$ gives $C(x_0)=C^\star(x_0)$
and strictly decreases $\mathcal{L}_{\mathrm{div}}(x_0)$ compared to the
collapsed configuration.

\medskip
\noindent\textbf{Step 2: Segmentation reconstruction under collapse.}
By assumption (ii), the segmentation loss $\mathcal{L}_{\mathrm{seg}}$
is minimized only when $S(x) = S^\star(x)$ for every image $x$.
Under the collapsed configuration for $x_0$, we have $S(x_0)\equiv c$.
Since $C^\star(x_0) \neq 0$, the teacher embeddings $S^\star(x_0)$
cannot all be equal to the same constant vector, so $S(x_0) \neq S^\star(x_0)$.
Thus replacing $S(x_0)$ by $S^\star(x_0)$ strictly decreases
$\mathcal{L}_{\mathrm{seg}}$.

\medskip
\noindent\textbf{Step 3: Alignment under collapse.}
By assumption (iii), $\mathcal{L}_{\mathrm{align}}$ is minimized only when
the collection $\{(G(x), S(x))\}_x$ is discriminative across images.
In particular, a configuration in which mask embeddings are constant
or nearly constant across images does not minimize $\mathcal{L}_{\mathrm{align}}$.
Therefore, there exists a perturbation of the collapsed configuration
that strictly decreases $\mathcal{L}_{\mathrm{align}}$.

\medskip
\noindent\textbf{Step 4: Contradiction.}
Starting from the collapsed configuration for $x_0$, consider a modified
configuration in which we set $S(x_0) = S^\star(x_0)$ and (if necessary)
adjust parameters to reduce $\mathcal{L}_{\mathrm{align}}$, while keeping
all other images fixed.
From Steps 1 and 2, this modification strictly decreases
$\mathcal{L}_{\mathrm{div}}(x_0)$ and $\mathcal{L}_{\mathrm{seg}}$.
From Step 3, we can also strictly decrease $\mathcal{L}_{\mathrm{align}}$.
Since $\lambda_3,\lambda_4>0$, any strict decrease in one of
$\mathcal{L}_{\mathrm{seg}}, \mathcal{L}_{\mathrm{align}},
\mathcal{L}_{\mathrm{div}}$ leads to a strict decrease in
\[
\mathcal{L}_{\mathrm{emb}}
= \mathcal{L}_{\mathrm{seg}}
+ \lambda_3 \mathcal{L}_{\mathrm{align}}
+ \lambda_4 \mathcal{L}_{\mathrm{div}}.
\]
Thus the collapsed configuration cannot be a global minimizer of
$\mathcal{L}_{\mathrm{emb}}$, contradicting our assumption.

Therefore, no global minimizer of $\mathcal{L}_{\mathrm{emb}}$ can have
all predicted mask embeddings collapse to a single constant vector
for any image $x$.
\end{proof}

\subsection{Structural Sufficiency of I-JEPA + SAM Representations}

In this section we formalize a restricted sufficiency property of the
joint representation \(Z = (G,S)\), where \(G\) is the global I-JEPA
embedding and \(S\) is the SAM-based per-mask embedding. This property
justifies using \(Z\) as the latent variable in our generative
factorization.

\subsubsection{Setup}

Let \(X\) be a random variable of images with distribution \(p_X\).
Let \(E_{\mathrm{g}} : \mathcal{X} \to \mathbb{R}^{d_g}\) denote the
I-JEPA encoder, and \(E_{\mathrm{s}} : \mathcal{X} \to \mathbb{R}^{N \times d_s}\)
denote the SAM-based encoder that outputs per-mask embeddings. We define
\[
G = E_{\mathrm{g}}(X) \in \mathbb{R}^{d_g}, \qquad
S = E_{\mathrm{s}}(X) \in \mathbb{R}^{N \times d_s}, \qquad
Z = (G,S).
\]

Let \(\mathcal{R}\) be the space of region descriptors (e.g., indices or
positional encodings of masked regions). For each image \(x \in \mathcal{X}\)
and region descriptor \(R \in \mathcal{R}\), let
\[
Y_R(x) \in \mathbb{R}^{d_y}
\]
denote the corresponding teacher region embedding, as defined by the
I-JEPA teacher network.\footnote{For example, \(Y_R(x)\) may be the
teacher ViT embedding associated with the patch subset specified by \(R\)
in the formulation of Assran et al.~\cite{assran2023ijepa}.}

The I-JEPA training objective for the image-level encoder \(E_{\mathrm{g}}\)
and predictor \(P\) can be written (in squared-error form) as
\begin{equation}
\label{eq:ijepa-objective}
\min_{\theta,\psi} \;
\mathbb{E}_{X \sim p_X}\,\mathbb{E}_{R \sim \mathcal{R}}
\big[
\|P_\psi(E_{\mathrm{g},\theta}(X), R) - Y_R(X)\|_2^2
\big],
\end{equation}
where \(\theta,\psi\) are the parameters of \(E_{\mathrm{g}}\) and \(P\),
respectively.

\subsubsection{Structural Sufficiency of \texorpdfstring{$(G,S)$}{(G,S)}}

We now show that, under standard regression assumptions, the global
I-JEPA embedding \(G\) is a sufficient statistic (for this family of
region-level prediction tasks), and hence the joint representation
\(Z = (G,S)\) is also sufficient.

\begin{proposition}[Structural sufficiency of I-JEPA + SAM representation]
\label{prop:structural-sufficiency}
Assume that:
\begin{enumerate}
    \item[(i)] \(E_{\mathrm{g}}\) and \(P\) are trained by minimizing
    the objective \eqref{eq:ijepa-objective} with squared error,
    and that the teacher region embeddings \(Y_R(X)\) are fixed
    (non-trainable) random variables;
    \item[(ii)] the predictor class is rich enough that the empirical
    minimizer \(P^\star\) of \eqref{eq:ijepa-objective} attains the
    regression optimum for almost all \((G,R)\).
\end{enumerate}
Then for almost all \((g,R)\), the optimal predictor satisfies
\begin{equation}
\label{eq:cond-mean-g}
P^\star(g,R)
= \mathbb{E}\big[\,Y_R(X) \mid G = g, R\,\big].
\end{equation}
Moreover, since \(G = E_{\mathrm{g}}(X)\) and
\(Z = (G,S) = (E_{\mathrm{g}}(X), E_{\mathrm{s}}(X))\) are deterministic
functions of \(X\), we also have
\begin{equation}
\label{eq:cond-mean-z}
\mathbb{E}\big[\,Y_R(X) \mid G, R\,\big]
= \mathbb{E}\big[\,Y_R(X) \mid Z, R\,\big],
\end{equation}
so the joint representation \(Z = (G,S)\) is a sufficient statistic for
predicting the family of targets \(\{Y_R(X)\}_{R \in \mathcal{R}}\) in the
sense that all optimal predictors of \(Y_R(X)\) given \(X\) and \(R\) can
be expressed as functions of \((Z,R)\).
\end{proposition}

\begin{proof}
We first prove \eqref{eq:cond-mean-g}. For any fixed encoder
\(E_{\mathrm{g}}\), the optimization problem \eqref{eq:ijepa-objective}
over predictors \(P\) is, for each region descriptor \(R\), a standard
least-squares regression of \(Y_R(X)\) on \(G = E_{\mathrm{g}}(X)\) and
\(R\). Under assumptions (i)–(ii), the minimizer \(P^\star\) of the
expected squared loss satisfies the well-known characterization of
conditional expectation in \(L^2\):
\[
P^\star(G,R)
= \mathbb{E}\big[\,Y_R(X) \mid G, R\,\big]
\quad\text{almost surely},
\]
which yields \eqref{eq:cond-mean-g}.

Next, we establish \eqref{eq:cond-mean-z}. Recall that both \(G\) and
\(Z = (G,S)\) are deterministic functions of the underlying image \(X\):
\[
G = E_{\mathrm{g}}(X), \qquad
Z = (E_{\mathrm{g}}(X), E_{\mathrm{s}}(X)).
\]
Therefore, for any measurable function \(h\),
\[
h(X) = h(X, G, Z)
\]
can be viewed as a function of \(X\) alone or as a function of
\((X,G,Z)\), and conditioning on \((G,R)\) or on \((Z,R)\) induces
the same \(\sigma\)-algebra generated by \(G\) together with \(R\), up to
augmentation by the additional information in \(S\) that is itself a
measurable function of \(X\).

Formally, because \(G = E_{\mathrm{g}}(X)\) and \(Z = (G,S)\) with
\(S = E_{\mathrm{s}}(X)\), we have the inclusion of \(\sigma\)-algebras
\[
\sigma(G,R) \subseteq \sigma(Z,R).
\]
Thus
\[
\mathbb{E}\big[\,Y_R(X) \mid Z, R\,\big]
= \mathbb{E}\big[
\mathbb{E}\big[\,Y_R(X) \mid G, R\,\big] \,\big|\, Z, R
\big]
= \mathbb{E}\big[\,Y_R(X) \mid G, R\,\big],
\]
where the first equality uses the tower property of conditional
expectation and the fact that \(\mathbb{E}[Y_R(X)\mid G,R]\) is already
\(\sigma(G,R)\)-measurable and hence also \(\sigma(Z,R)\)-measurable.
This proves \eqref{eq:cond-mean-z}.

Combining \eqref{eq:cond-mean-g} and \eqref{eq:cond-mean-z}, we conclude
that the optimal predictor of \(Y_R(X)\) can be written as a function of
\((G,R)\), and equivalently as a function of \((Z,R)\).
In particular, any optimal predictor of the teacher’s region embeddings
can be implemented using only the joint representation \(Z = (G,S)\) and
the region descriptor \(R\).
This establishes the claimed sufficiency of \(Z\) for the family of
prediction tasks \(\{Y_R(X)\}_{R \in \mathcal{R}}\).
\end{proof}

\subsection{Partition-Based Spatial Modulation with SAM}

In this section we formalize how SAM's instance masks and per-mask
embeddings induce a natural, minimal spatially adaptive conditioning
field for our segmentation-aware generator.

\subsubsection{Setup}

Let the image domain be a discrete grid
\(\Omega = \{1,\dots,H\} \times \{1,\dots,W\}\).
For a fixed image \(x\), SAM produces a set of non-overlapping
instance masks
\[
\{M_i(x)\}_{i=1}^N, \quad M_i(x) \subset \Omega,
\]
such that
\[
M_i(x) \cap M_j(x) = \emptyset \ \text{for } i \neq j,
\qquad
\bigcup_{i=1}^N M_i(x) = \Omega.
\]
Thus \(\{M_i(x)\}_{i=1}^N\) forms a partition of the image grid.
For each region \(M_i(x)\) we have a per-mask embedding
\(s_i(x) \in \mathbb{R}^{d_s}\) provided by the SAM encoder.

Let \(C\) be the channel dimension of a feature map
\(x^{\mathrm{feat}} \in \mathbb{R}^{C \times H \times W}\) inside the
generator.
We consider spatially adaptive modulation parameters
\(\gamma(x),\beta(x) \in \mathbb{R}^{C \times H \times W}\) that are
piecewise constant on the SAM partition and are functions only of the
per-mask embeddings \(\{s_i(x)\}\).

\subsubsection{Minimal Partition-Based Conditioner}

We now show that, given the SAM partition and per-mask embeddings, our
construction of \((\gamma(x),\beta(x))\) is the unique piecewise-constant
field adapted to the partition that depends only on region-level codes.

\begin{lemma}[Partition-based spatial modulation]
\label{lem:partition-modulation}
Fix an image \(x\) and its SAM partition
\(\{M_i(x)\}_{i=1}^N\) of \(\Omega\).
Let \(s_i(x) \in \mathbb{R}^{d_s}\) be the per-region codes and let
\(W_\gamma, W_\beta : \mathbb{R}^{d_s} \to \mathbb{R}^C\) be linear maps.
Define region-level modulation parameters
\[
\gamma_i(x) = W_\gamma s_i(x) \in \mathbb{R}^C, \qquad
\beta_i(x) = W_\beta s_i(x) \in \mathbb{R}^C.
\]
For each spatial location \((u,v) \in \Omega\), let \(i(u,v)\) be the
unique index such that \((u,v) \in M_{i(u,v)}(x)\).
Define the spatial fields \(\gamma(x),\beta(x) \in \mathbb{R}^{C \times H \times W}\) by
\[
\gamma_c(u,v) = \gamma_{i(u,v),c}(x), \qquad
\beta_c(u,v) = \beta_{i(u,v),c}(x),
\quad c = 1,\dots,C.
\]

Then:
\begin{enumerate}
    \item[(a)] \(\gamma(x)\) and \(\beta(x)\) are piecewise constant on
    the partition \(\{M_i(x)\}\), i.e., for each \(i\), the restrictions
    of \(\gamma(x)\) and \(\beta(x)\) to \(M_i(x)\) are constant;
    \item[(b)] conversely, any spatial fields
    \(\tilde{\gamma}(x),\tilde{\beta}(x) \in \mathbb{R}^{C \times H \times W}\)
    that are piecewise constant on \(\{M_i(x)\}\) and depend only on the
    region-level codes \(\{s_i(x)\}_{i=1}^N\) can be written in the form
    above for some (possibly different) linear maps
    \(W_\gamma, W_\beta : \mathbb{R}^{d_s} \to \mathbb{R}^C\).
\end{enumerate}
In particular, SAM's partition and per-mask embeddings determine a
minimal spatially adaptive conditioning mechanism in which each region
\(M_i(x)\) is assigned a constant modulation vector that depends only on
its code \(s_i(x)\).
\end{lemma}

\begin{proof}
(a) By construction, for any fixed region index \(i\) and any two
locations \((u,v), (u',v') \in M_i(x)\), we have
\(i(u,v) = i(u',v') = i\).
Thus
\[
\gamma_c(u,v) = \gamma_{i,c}(x), \quad
\gamma_c(u',v') = \gamma_{i,c}(x),
\]
and similarly for \(\beta\).
Hence \(\gamma(x)\) and \(\beta(x)\) are constant on each region
\(M_i(x)\), and therefore piecewise constant on the partition
\(\{M_i(x)\}\).

(b) Suppose \(\tilde{\gamma}(x),\tilde{\beta}(x)\) are spatial fields that
are piecewise constant on \(\{M_i(x)\}\) and depend only on
\(\{s_i(x)\}_{i=1}^N\).
Being piecewise constant means that for each \(i\) there exist vectors
\(\tilde{\gamma}_i(x), \tilde{\beta}_i(x) \in \mathbb{R}^C\) such that
\[
\tilde{\gamma}_c(u,v) = \tilde{\gamma}_{i,c}(x), \qquad
\tilde{\beta}_c(u,v) = \tilde{\beta}_{i,c}(x),
\quad \text{for all } (u,v) \in M_i(x), \ c=1,\dots,C.
\]
Dependence only on the region codes \(\{s_i(x)\}\) means that
\(\tilde{\gamma}_i(x)\) and \(\tilde{\beta}_i(x)\) can be written as
functions of \(s_i(x)\) alone:
\[
\tilde{\gamma}_i(x) = F_\gamma(s_i(x)), \qquad
\tilde{\beta}_i(x) = F_\beta(s_i(x)),
\]
for some functions \(F_\gamma, F_\beta : \mathbb{R}^{d_s} \to \mathbb{R}^C\).

If we restrict to linear dependence on \(s_i(x)\), then
\(F_\gamma, F_\beta\) are linear maps and can be represented as
\(W_\gamma, W_\beta \in \mathbb{R}^{C \times d_s}\).
In that case, \(\tilde{\gamma}(x),\tilde{\beta}(x)\) coincide with the
fields defined in the statement of the lemma:
\[
\tilde{\gamma}_i(x) = W_\gamma s_i(x), \qquad
\tilde{\beta}_i(x) = W_\beta s_i(x),
\]
and
\[
\tilde{\gamma}_c(u,v) = \tilde{\gamma}_{i(u,v),c}(x), \quad
\tilde{\beta}_c(u,v) = \tilde{\beta}_{i(u,v),c}(x).
\]

Therefore, within the class of linear region-level conditioners, any
piecewise-constant spatial modulation field adapted to the SAM partition
and depending only on \(\{s_i(x)\}\) necessarily has the form described
in the lemma.
\end{proof}

\begin{remark}[Cross-attention as a soft partition conditioner]
In our implementation, we replace the hard partition-based assignment
\((u,v) \mapsto i(u,v)\) with a learned cross-attention mechanism:
at each spatial location the feature map queries all segment tokens
\(\{s_i(x)\}\) via multi-head attention, yielding a soft, data-dependent
weighting.  This is strictly more expressive than the piecewise-constant
field described in Lemma~\ref{lem:partition-modulation}: by concentrating
attention weights on a single token per location, the cross-attention layer
can recover the partition-based modulation as a special case, while also
allowing smooth blending across region boundaries.
A gated residual connection
\[
x \leftarrow x + \alpha \cdot \mathrm{CrossAttn}(x,\, S),
\]
with \(\alpha\) initialised at zero, ensures that training begins from
the unconditional baseline and gradually incorporates structural
conditioning as the gate opens.
\end{remark}

%%%%%%%%%%%%%%%%%%%%%%%%%%%%%%%%%%%%%%%%%%%%%%%%%%%%%%%%%%%%%%%%%%%%%%%%%%%%%%%%%%%%%%%%%%%%%%%%%%%%%%%%%%%%%%%%%%%%%%%%%%%%%%%%%%%%%%%%%%%%%
\section{Problem Setup and Notation}
\label{sec:setup}
%% ====================================================================

Let $\mathcal{X} \subseteq \mathbb{R}^{3 \times H \times W}$ denote the
space of RGB images of resolution $H \times W$, and let
$X \sim p_{\mathrm{data}}$ be a random image drawn from the data distribution.
We employ two frozen, pretrained encoders:
\begin{enumerate}
    \item \textbf{I-JEPA encoder}
    (Image-based Joint-Embedding Predictive Architecture~\cite{assran2023ijepa}):
    a ViT-H/14 context encoder
    $\phi_{\mathrm{jepa}} : \mathcal{X} \to \mathbb{R}^{d_g}$
    that maps an image to a global semantic representation $g = \phi_{\mathrm{jepa}}(x)$
    with $d_g = 1280$.
    In practice, the raw ViT output is projected to $\mathbb{R}^{256}$ via
    a learned linear projection $W_{\mathrm{proj}} \in \mathbb{R}^{256 \times 1280}$.
    We write $g \in \mathbb{R}^{256}$ for the projected embedding.
    \item \textbf{SAM encoder}
    (Segment Anything Model~\cite{kirillov2023sam}):
    an image encoder followed by per-mask average pooling,
    $\phi_{\mathrm{sam}} : \mathcal{X} \to \mathbb{R}^{N \times d_s}$,
    that yields $N$ segment-level embeddings
    $s_1, \ldots, s_N \in \mathbb{R}^{d_s}$
    with $d_s = 256$.
    Each $s_i$ is the spatially pooled SAM feature map restricted to
    the $i$-th mask region $M_i \subset \Omega$,
    where $\{M_i\}_{i=1}^{N}$ is a partition of the image grid
    $\Omega = \{1,\ldots,H\}\times\{1,\ldots,W\}$.
\end{enumerate}

For a given image $x$, we define the \emph{unified token sequence}
\begin{equation}\label{eq:token-sequence}
    z = \bigl[\, g,\; s_1,\; s_2,\; \ldots,\; s_N \,\bigr]
    \;\in\; \mathbb{R}^{(1+N) \times 256},
\end{equation}
where the first token is the I-JEPA global vector and the remaining
$N$ tokens are SAM segment embeddings.
In practice, sequences are padded to a fixed length $L = 251$ with
zero vectors and a boolean padding mask $m \in \{0,1\}^L$ is maintained.

%% ====================================================================
\section{Two-Stage Generative Model}
\label{sec:two-stage}
%% ====================================================================

Our generative framework factorises image synthesis into two stages.

\begin{definition}[Two-stage factorisation]
\label{def:two-stage}
Define the latent variable $z = (g, S)$ with $g \in \mathbb{R}^{256}$ and
$S = (s_1, \ldots, s_N) \in \mathbb{R}^{N \times 256}$.
The data distribution is decomposed as
\begin{equation}\label{eq:marginalisation}
    p_{\mathrm{data}}(x)
    = \int p(x \mid g, S)\, p(g, S)\, dg\, dS,
\end{equation}
and our model approximates each factor:
\begin{itemize}
    \item \textbf{Stage~1} (Representation Diffusion Model):
    a denoising diffusion probabilistic model (DDPM)~\cite{ho2020ddpm}
    that learns $p_\theta(g, S) \approx p(g, S)$
    in the joint representation space.
    \item \textbf{Stage~2} (Structure-Aware Generator):
    a conditional StyleGAN2-based generator~\cite{karras2020stylegan2}
    that learns $p_\psi(x \mid g, S) \approx p(x \mid g, S)$
    via adversarial training with structural conditioning.
\end{itemize}
\end{definition}

At inference, we sample
$\hat{z} = (\hat{g}, \hat{S}) \sim p_\theta$
and decode
$\hat{x} = G_\psi(\hat{g}, \hat{S})$,
yielding an overall generative model
\begin{equation}\label{eq:generative-model}
    \hat{p}(x) = \int p_\psi(x \mid g, S)\, p_\theta(g, S)\, dg\, dS.
\end{equation}

%% ====================================================================
\section{Stage~1: Denoising Diffusion in Representation Space}
\label{sec:stage1}
%% ====================================================================

\subsection{Forward Diffusion Process}

We adopt the standard DDPM framework~\cite{ho2020ddpm,sohl2015deep}
operating in the 256-dimensional representation space
rather than pixel space.
Let $z_0 \in \mathbb{R}^{L \times 256}$ denote the clean unified
token sequence for a training example.
The forward (noising) process is a fixed Markov chain
$q(z_t \mid z_{t-1}) = \mathcal{N}(z_t; \sqrt{1-\beta_t}\, z_{t-1},\, \beta_t I)$
for $t = 1, \ldots, T$ with $T = 1000$.

Defining $\alpha_t = 1 - \beta_t$ and $\bar{\alpha}_t = \prod_{k=1}^{t} \alpha_k$,
the marginal at arbitrary timestep $t$ admits the closed form
\begin{equation}\label{eq:forward-marginal}
    q(z_t \mid z_0)
    = \mathcal{N}\!\bigl(z_t;\; \sqrt{\bar{\alpha}_t}\, z_0,\;
    (1 - \bar{\alpha}_t)\, I\bigr),
\end{equation}
equivalently
$z_t = \sqrt{\bar{\alpha}_t}\, z_0 + \sqrt{1-\bar{\alpha}_t}\, \epsilon$
with $\epsilon \sim \mathcal{N}(0, I)$.

\paragraph{Noise schedule.}
We use a \emph{linear} schedule~\cite{ho2020ddpm}:
$\beta_t = \bigl(\sqrt{\beta_1} + \frac{t-1}{T-1}\bigl(\sqrt{\beta_T} - \sqrt{\beta_1}\bigr)\bigr)^{\!2}$
with $\beta_1 = 10^{-4}$ and $\beta_T = 0.02$.

\subsection{Denoising Objective ($x_0$-Prediction)}
\label{sec:x0-pred}

The reverse process
$p_\theta(z_{t-1} \mid z_t) = \mathcal{N}(z_{t-1}; \mu_\theta(z_t, t), \sigma_t^2 I)$
is parameterised via \emph{$x_0$-prediction}~\cite{ramesh2022hierarchical}:
the denoising network $f_\theta(z_t, t)$ directly predicts $\hat{z}_0$,
and the posterior mean is obtained from the DDPM posterior
\begin{equation}\label{eq:posterior-mean}
    \mu_\theta(z_t, t)
    = \frac{\sqrt{\bar{\alpha}_{t-1}}\,\beta_t}{1 - \bar{\alpha}_t}\,
    f_\theta(z_t, t)
    \;+\;
    \frac{\sqrt{\alpha_t}\,(1 - \bar{\alpha}_{t-1})}{1 - \bar{\alpha}_t}\,
    z_t.
\end{equation}

The training objective is the \emph{simple loss}~\cite{ho2020ddpm}:
\begin{equation}\label{eq:ddpm-loss}
    \mathcal{L}_{\mathrm{simple}}(\theta)
    = \mathbb{E}_{z_0,\, t,\, \epsilon}\!
    \bigl[\,\lVert f_\theta(z_t, t) - z_0 \rVert_2^2 \,\bigr],
    \quad t \sim \mathrm{Uniform}\{1, \ldots, T\}.
\end{equation}

When the token sequence contains padded positions (indicated by the mask $m$),
the loss is computed only over valid (non-padded) tokens:
\begin{equation}\label{eq:masked-loss}
    \mathcal{L}_{\mathrm{simple}}(\theta)
    = \mathbb{E}_{z_0,\, t,\, \epsilon}\!\biggl[\,
    \frac{1}{\sum_{\ell} m_\ell}\,
    \sum_{\ell=1}^{L} m_\ell\,
    \lVert f_\theta(z_t, t)_\ell - (z_0)_\ell \rVert_2^2
    \,\biggr].
\end{equation}

\begin{remark}[$x_0$-prediction vs.\ $\epsilon$-prediction]
\label{rem:x0-vs-eps}
Standard DDPM~\cite{ho2020ddpm} trains the network to predict the added noise
$\epsilon$.  In the $x_0$-prediction parameterisation
(also used in DALL$\cdot$E~2~\cite{ramesh2022hierarchical} and
RDM~\cite{bordes2022rdm}),
the network instead predicts the clean data $z_0$ directly.
The two are related by the identity
$\hat{\epsilon} = (z_t - \sqrt{\bar{\alpha}_t}\,\hat{z}_0) / \sqrt{1 - \bar{\alpha}_t}$.
The $x_0$ parameterisation is particularly natural for representation
diffusion since the target representations are bounded and low-dimensional.
\end{remark}

\subsection{Denoising Network Architecture}
\label{sec:denoising-arch}

We employ two backbone architectures depending on the operating regime.

\paragraph{SimpleMLP (single-vector RDM).}
For diffusing only the global I-JEPA vector $g \in \mathbb{R}^{256}$,
we use a residual MLP with $K = 12$ blocks,
hidden dimension $d_h = 1536$,
and sinusoidal timestep embeddings~\cite{vaswani2017attention}:
\begin{equation}\label{eq:mlp-block}
    h^{(k+1)} = h^{(k)} + \mathrm{MLP}\!\bigl(\mathrm{LN}(h^{(k)}) + e_t\bigr),
    \quad k = 0, \ldots, K-1,
\end{equation}
where $h^{(0)} = W_{\mathrm{in}}\, z_t$, $e_t = \mathrm{MLP}_t(\mathrm{SinEmb}(t))$,
and the final output is $\hat{z}_0 = W_{\mathrm{out}}\, \mathrm{LN}(h^{(K)})$
with $W_{\mathrm{out}}$ zero-initialized~\cite{goyal2017accurate}.
\subsection{Auxiliary Losses for Representation Quality}
\label{sec:aux-losses}

We augment the simple denoising loss~\eqref{eq:ddpm-loss} with two
auxiliary terms that operate on the \emph{predicted clean tokens}
$\hat{z}_0 = f_\theta(z_t, t)$.
Let $\hat{g} = (\hat{z}_0)_0 \in \mathbb{R}^{256}$ denote the predicted
global vector and $\hat{S} = (\hat{s}_1, \ldots, \hat{s}_N)$ the predicted
segment tokens.

\paragraph{Diversity loss (log-determinant regulariser).}
Define the empirical covariance of the predicted segment tokens:
\begin{equation}\label{eq:seg-cov}
    \bar{s} = \frac{1}{N}\sum_{i=1}^{N} \hat{s}_i,
    \qquad
    C = \frac{1}{N}\sum_{i=1}^{N}
    (\hat{s}_i - \bar{s})(\hat{s}_i - \bar{s})^\top
    \in \mathbb{R}^{256 \times 256}.
\end{equation}
The diversity loss is
\begin{equation}\label{eq:div-loss}
    \mathcal{L}_{\mathrm{div}} = -\log\det(C + \epsilon\, I),
    \quad \epsilon = 10^{-4},
\end{equation}
which penalises low-rank covariance (i.e., collapsed or redundant
segment embeddings).
This is equivalent to the negative log-volume of the parallelepiped
spanned by the centred embeddings, and is related to Determinantal
Point Process (DPP) kernels~\cite{kulesza2012dpp} and the VICReg
variance term~\cite{bardes2022vicreg}.

\paragraph{Alignment loss (InfoNCE).}
We use an InfoNCE contrastive loss~\cite{oord2018infonce} between
the predicted global vector and segment tokens within each mini-batch
of size $B$:
\begin{equation}\label{eq:align-loss}
    \mathcal{L}_{\mathrm{align}}
    = -\frac{1}{B}\sum_{i=1}^{B}
    \log \frac{
    \exp\!\bigl(\mathrm{sim}(\hat{g}_i,\, \bar{s}_i)\,/\,\tau\bigr)
    }{
    \sum_{j=1}^{B}
    \exp\!\bigl(\mathrm{sim}(\hat{g}_i,\, \bar{s}_j)\,/\,\tau\bigr)
    },
\end{equation}
where $\mathrm{sim}(a, b) = a^\top b / (\lVert a \rVert \lVert b \rVert)$
is cosine similarity and $\tau = 0.07$ is the temperature.
This encourages the global and segment representations to remain
mutually informative and image-specific.

\paragraph{Total Stage-1 loss.}
\begin{equation}\label{eq:total-stage1-loss}
    \mathcal{L}_{\mathrm{Stage1}}
    = \mathcal{L}_{\mathrm{simple}}
    + \lambda_{\mathrm{div}}\, \mathcal{L}_{\mathrm{div}}
    + \lambda_{\mathrm{align}}\, \mathcal{L}_{\mathrm{align}},
\end{equation}
with $\lambda_{\mathrm{div}} = 0.1$ and $\lambda_{\mathrm{align}} = 0.05$.

\subsection{Sampling via DDIM}
\label{sec:ddim}

At inference, we use the DDIM sampler~\cite{song2021ddim}
for accelerated, optionally deterministic generation.
Given a subsequence of timesteps
$\{t_1 > t_2 > \cdots > t_S\} \subset \{1, \ldots, T\}$
with $S \ll T$ (typically $S = 50$), the DDIM update is:
\begin{equation}\label{eq:ddim-update}
    z_{t_{i+1}} = \sqrt{\bar{\alpha}_{t_{i+1}}}\,\hat{z}_0
    + \sqrt{1 - \bar{\alpha}_{t_{i+1}} - \sigma_{t_i}^2}\;\hat{\epsilon}
    + \sigma_{t_i}\, \xi,
    \quad \xi \sim \mathcal{N}(0, I),
\end{equation}
where $\hat{z}_0 = f_\theta(z_{t_i}, t_i)$ (the $x_0$-prediction),
$\hat{\epsilon} = (z_{t_i} - \sqrt{\bar{\alpha}_{t_i}}\,\hat{z}_0) / \sqrt{1 - \bar{\alpha}_{t_i}}$
(the implied noise prediction), and
\begin{equation}\label{eq:ddim-sigma}
    \sigma_{t_i} = \eta\,\sqrt{\frac{1 - \bar{\alpha}_{t_{i+1}}}{1 - \bar{\alpha}_{t_i}}
    \cdot \frac{1 - \bar{\alpha}_{t_i}/\bar{\alpha}_{t_{i+1}}}{1}}.
\end{equation}
Setting $\eta = 0$ yields deterministic sampling; $\eta = 1$ recovers
DDPM ancestral sampling.

%% ====================================================================
\section{Stage~2: Structure-Aware Adversarial Decoder}
\label{sec:stage2}
%% ====================================================================

Stage~2 is a conditional GAN that decodes sampled representations
$(g, S)$ into images.
The generator and discriminator are built upon
StyleGAN2-ADA~\cite{karras2020stylegan2,karras2020training}.

\subsection{Generator Architecture}

\subsubsection{I-JEPA Conditioning via $\mathcal{W}$-Space Fusion}

The I-JEPA global vector $g \in \mathbb{R}^{256}$ is injected into the
StyleGAN2 generator through the $\mathcal{W}$ latent space.
A learnable projection network
$\pi : \mathbb{R}^{256} \to \mathbb{R}^{w}$
(two fully-connected layers with Leaky ReLU, where $w = 512$ is the
$\mathcal{W}$-space dimension) maps $g$ to a style modifier.
Given the standard StyleGAN2 mapping network output
$\mathbf{w} = (\mathbf{w}_1, \ldots, \mathbf{w}_{L_w}) \in \mathbb{R}^{L_w \times w}$,
we apply an \emph{additive semantic shift} to the first $K$ layers
($K = 6$ by default):
\begin{equation}\label{eq:w-fusion}
    \tilde{\mathbf{w}}_k =
    \begin{cases}
    \mathbf{w}_k + \alpha_{\mathrm{sem}} \cdot r_{\mathrm{sem}}
    \cdot \mathrm{norm}(\pi(g)),
    & k \leq K, \\
    \mathbf{w}_k, & k > K,
    \end{cases}
\end{equation}
where $\alpha_{\mathrm{sem}} = \tanh(\alpha_{\mathrm{raw}})$ is a learnable
gating scalar (initialised at $\alpha_{\mathrm{raw}} = 0.3$),
$r_{\mathrm{sem}} \in [0,1]$ is a linear warmup ramp, and
$\mathrm{norm}(\cdot)$ denotes second-moment normalisation
as in~\cite{karras2020stylegan2}.

During training, a \emph{stochastic mixing depth} is sampled per
example: $K_i \sim \mathrm{Uniform}\{1,\ldots,K\}$, so that
layers beyond $K_i$ receive no semantic shift.
This acts as a regulariser analogous to style mixing~\cite{karras2019style}.

\subsubsection{SAM Conditioning via Cross-Attention}
\label{sec:cross-attn}

Structural information from SAM segment tokens is injected into the
synthesis network via \emph{cross-attention layers} inserted at
intermediate resolutions ($16 \times 16$ and $32 \times 32$).

\begin{definition}[Segment cross-attention]
\label{def:seg-cross-attn}
Let $x \in \mathbb{R}^{C \times H' \times W'}$ be a feature map at some
resolution $H' \times W'$ inside the synthesis network,
and let $S = (s_1, \ldots, s_N) \in \mathbb{R}^{N \times 256}$
be the SAM segment tokens.
\begin{enumerate}
    \item Flatten and project the feature map to queries:
    $Q = \mathrm{LN}(W_Q\, x_{\mathrm{flat}}) \in \mathbb{R}^{(H'W') \times d}$,
    where $x_{\mathrm{flat}} \in \mathbb{R}^{(H'W') \times C}$ and $d = 256$.
    \item Project segment tokens to keys and values:
    $K = V = \mathrm{LN}(W_S\, S) \in \mathbb{R}^{N \times d}$.
    \item Compute multi-head attention with $n_h = 8$ heads:
    $\mathrm{Attn}(Q, K, V) = \softmax\!\bigl(QK^\top / \sqrt{d/n_h}\bigr)\, V$,
    using the padding mask to exclude invalid tokens.
    \item Project back to the original channel dimension and apply a
    \emph{gated residual connection}:
    \begin{equation}\label{eq:gated-residual}
        x \leftarrow x + \alpha_{\mathrm{seg}} \cdot r_{\mathrm{seg}}
        \cdot W_O\, \mathrm{Attn}(Q, K, V),
    \end{equation}
    where $\alpha_{\mathrm{seg}}$ is a \emph{learnable scalar gate}
    initialised to zero, and $r_{\mathrm{seg}} \in [0,1]$ is a
    linear warmup ramp.
\end{enumerate}
\end{definition}

\begin{remark}[Zero-initialised gating]
\label{rem:zero-gate}
Initialising $\alpha_{\mathrm{seg}} = 0$ ensures that the generator
begins training as an \emph{unconditional} StyleGAN2 and
gradually incorporates structural conditioning as the gate opens.
This stabilises early training when the SAM embeddings from Stage~1
may not yet be reliable
(cf.\ zero-init residuals in~\cite{bachlechner2021rezero,goyal2017accurate}).
\end{remark}

\subsubsection{Stochastic SAM Dropout}

During training, segment conditioning is dropped with probability
$p_{\mathrm{seg}} = 0.15$ per batch element (setting $S = \mathbf{0}$
and disabling cross-attention).
This serves two purposes:
(i)~it regularises the generator to not rely exclusively on segment
tokens, and
(ii)~at inference it enables a form of classifier-free guidance
(CFG)~\cite{ho2022cfg} by interpolating between the conditioned and
unconditioned outputs.

\subsection{Discriminator Architecture}

The discriminator follows the standard StyleGAN2 residual architecture
and is augmented with \emph{projection-based conditioning}~\cite{miyato2018projection}
from both I-JEPA and SAM representations.

\paragraph{I-JEPA projection.}
The global vector $g$ is projected to a conditioning map
$c_g = W_g\, g \in \mathbb{R}^{d_{\mathrm{cmap}}}$ via a
two-layer MLP.

\paragraph{SAM segment pooling.}
A learnable query token $q_{\mathrm{pool}} \in \mathbb{R}^{256}$
attends to all SAM segment tokens via multi-head attention
(8~heads), producing a single pooled vector
$\bar{s}_{\mathrm{pool}} \in \mathbb{R}^{256}$, which is then projected
to $c_s = W_s\, \bar{s}_{\mathrm{pool}} \in \mathbb{R}^{d_{\mathrm{cmap}}}$.

\paragraph{Blended conditioning.}
The final conditioning vector is a normalised weighted sum:
\begin{equation}\label{eq:d-cmap}
    c = \mathrm{norm}\!\bigl(w_1\, c_{\mathrm{class}} + w_2\, c_g + w_3\, c_s\bigr),
\end{equation}
where $c_{\mathrm{class}}$ is the optional class embedding,
$(w_1, w_2, w_3)$ are ramp-dependent weights,
and $\mathrm{norm}(\cdot)$ denotes second-moment normalisation.
The discriminator output is
$D(x, g, S) = \langle \psi(x),\, c \rangle$,
where $\psi(x) \in \mathbb{R}^{d_{\mathrm{cmap}}}$ is the
discriminator's feature representation of the image.

\subsection{Training Objectives}
\label{sec:stage2-losses}

\paragraph{Adversarial loss.}
We use the non-saturating logistic GAN objective~\cite{goodfellow2014gan}:
\begin{align}
    \mathcal{L}_D &= \mathbb{E}_{x \sim p_{\mathrm{data}}}\!
    \bigl[\softplus\!\bigl(-D(x, g, S)\bigr)\bigr]
    + \mathbb{E}_{z \sim p_z}\!
    \bigl[\softplus\!\bigl(D(G_\psi(z, g, S), g, S)\bigr)\bigr],
    \label{eq:d-loss} \\
    \mathcal{L}_G^{\mathrm{adv}} &= \mathbb{E}_{z \sim p_z}\!
    \bigl[\softplus\!\bigl(-D(G_\psi(z, g, S), g, S)\bigr)\bigr],
    \label{eq:g-adv-loss}
\end{align}
where $\softplus(u) = \log(1 + e^u)$.

\paragraph{I-JEPA feature matching.}
To encourage semantic consistency, the generated image's I-JEPA
embedding is matched to the conditioning vector via cosine similarity:
\begin{equation}\label{eq:fm-loss}
    \mathcal{L}_G^{\mathrm{fm}}
    = 1 - \mathrm{sim}\!\bigl(
    \phi_{\mathrm{jepa}}(G_\psi(z, g, S)),\; g
    \bigr),
\end{equation}
with a warmup weight $\lambda_{\mathrm{fm}} \cdot r(t)$ that linearly
increases from $0$ to $\lambda_{\mathrm{fm}}$.

\paragraph{R1 gradient penalty.}
Following~\cite{mescheder2018training}, we apply
\begin{equation}\label{eq:r1}
    \mathcal{L}_D^{\mathrm{R1}}
    = \frac{\gamma}{2}\,
    \mathbb{E}_{x \sim p_{\mathrm{data}}}\!\bigl[
    \lVert \nabla_x D(x, g, S) \rVert^2 \bigr],
\end{equation}
with $\gamma = 10$.

\paragraph{Path length regularisation.}
The standard StyleGAN2 path length penalty~\cite{karras2020stylegan2}
is applied to the generator to encourage smooth latent space geometry.

\paragraph{Total Stage-2 losses.}
\begin{align}
    \mathcal{L}_G &= \mathcal{L}_G^{\mathrm{adv}}
    + \lambda_{\mathrm{fm}}\, r(t)\, \mathcal{L}_G^{\mathrm{fm}}
    + \lambda_{\mathrm{pl}}\, \mathcal{L}_G^{\mathrm{pl}},
    \label{eq:total-g-loss} \\
    \mathcal{L}_D &= \mathcal{L}_D^{\mathrm{adv}}
    + \mathcal{L}_D^{\mathrm{R1}}.
    \label{eq:total-d-loss}
\end{align}

\subsection{RDM-Mixed Training}
\label{sec:rdm-mixed}

After a warmup period, Stage~2 training begins mixing in
\emph{synthetically sampled} representations from the Stage~1 diffusion
model.
With probability $p_{\mathrm{rdm}}$, the real embeddings $(g, S)$
for a training image are replaced with samples
$(\hat{g}, \hat{S}) \sim p_\theta$ drawn from the pretrained
Stage~1 model via DDIM ($S = 50$ steps, $\eta = 0$).
This exposes the generator and discriminator to the same distribution
they will encounter at inference, reducing the train--test mismatch
between real and generated representations.

%% ====================================================================
\section{Theoretical Analysis}
\label{sec:theory}
%% ====================================================================

\subsection{Diversity Loss Prevents Embedding Collapse}

\begin{lemma}[Diversity loss penalises collapsed embeddings]
\label{lem:div-collapse}
Let $S = (s_1, \ldots, s_N) \in \mathbb{R}^{N \times d}$ be a set of
segment embeddings with empirical covariance
$C = \frac{1}{N}\sum_{i=1}^{N}(s_i - \bar{s})(s_i - \bar{s})^\top$,
and let $\mathcal{L}_{\mathrm{div}} = -\log\det(C + \epsilon I)$
for $\epsilon > 0$.
Denote the eigenvalues of $C$ by $\lambda_1, \ldots, \lambda_d \geq 0$.
Then:
\begin{enumerate}
    \item[(a)] $\mathcal{L}_{\mathrm{div}} = -\sum_{k=1}^{d} \log(\lambda_k + \epsilon)$.
    \item[(b)] If all embeddings collapse to a single point
    ($s_i = c$ for all $i$), then $C = 0$ and
    $\mathcal{L}_{\mathrm{div}} = -d\log\epsilon$.
    \item[(c)] $\frac{\partial \mathcal{L}_{\mathrm{div}}}{\partial \lambda_k}
    = -(\lambda_k + \epsilon)^{-1} < 0$ for all $k$,
    so the loss strictly decreases as any eigenvalue increases from zero.
\end{enumerate}
\end{lemma}

\begin{proof}
Part (a) follows from the spectral decomposition
$C + \epsilon I = U \diag(\lambda_1 + \epsilon, \ldots, \lambda_d + \epsilon) U^\top$,
so $\det(C + \epsilon I) = \prod_{k}(\lambda_k + \epsilon)$ and
$\mathcal{L}_{\mathrm{div}} = -\sum_{k}\log(\lambda_k + \epsilon)$.

Part (b):\ if $s_i = c$ for all $i$, then $\bar{s} = c$ and $C = 0$,
giving $\det(\epsilon I) = \epsilon^d$ and
$\mathcal{L}_{\mathrm{div}} = -d\log\epsilon$.
Since $\epsilon \ll 1$, this is a large positive value.

Part (c):\ $-\log(\lambda_k + \epsilon)$ is strictly decreasing in
$\lambda_k$ for $\lambda_k \geq 0$.
\end{proof}

\subsection{InfoNCE Alignment Prevents Trivial Collapse}

\begin{lemma}[InfoNCE is maximised by trivial representations]
\label{lem:align-collapse}
Consider the alignment loss~\eqref{eq:align-loss} with
a batch $\{(\hat{g}_i, \hat{S}_i)\}_{i=1}^{B}$.
If all segment embeddings collapse to a constant
$\bar{s}_i = c$ for all $i$, then
$\mathrm{sim}(\hat{g}_i, \bar{s}_j) = \mathrm{sim}(\hat{g}_i, c)$
is independent of $j$, and
\begin{equation}
    \mathcal{L}_{\mathrm{align}} = \log B,
\end{equation}
which is the maximum of the InfoNCE loss.
Thus, minimising $\mathcal{L}_{\mathrm{align}}$ requires the
segment embeddings to carry image-specific information that
distinguishes positive from negative pairs.
\end{lemma}

\begin{proof}
Under collapse $\bar{s}_j = c$, the similarity
$\mathrm{sim}(\hat{g}_i, c)$ is the same for all $j$.
Hence the softmax in~\eqref{eq:align-loss} becomes uniform:
$\exp(\mathrm{sim}(\hat{g}_i, c)/\tau) / \sum_{j=1}^{B}\exp(\mathrm{sim}(\hat{g}_i, c)/\tau) = 1/B$
and $\mathcal{L}_{\mathrm{align}} = -\frac{1}{B}\sum_i \log(1/B) = \log B$.
\end{proof}

\subsection{Joint Loss Excludes Collapsed Minimisers}

\begin{proposition}[No collapsed global minimiser]
\label{prop:no-collapse}
Let $S^\star(x) = (s_1^\star(x), \ldots, s_N^\star(x))$ denote
the teacher SAM embeddings for image $x$.
Assume:
\begin{enumerate}
    \item[\textup{(A1)}] There exists at least one image $x$ for which the
    teacher covariance $C^\star(x) \neq 0$
    (i.e., the SAM embeddings are not degenerate).
    \item[\textup{(A2)}] The denoising loss $\mathcal{L}_{\mathrm{simple}}$
    is minimised only when the predicted tokens match the teacher tokens
    for every training example.
    \item[\textup{(A3)}] The alignment loss $\mathcal{L}_{\mathrm{align}}$
    is minimised only when $\{(\hat{g}_i, \hat{S}_i)\}_{i=1}^{B}$
    are discriminative across images.
\end{enumerate}
Then any parameter configuration in which the predicted segment
embeddings collapse to a single constant vector for some image $x$
(i.e., $\hat{s}_i(x) = c$ for all $i$) cannot be a global minimiser of
\begin{equation}
    \mathcal{L}_{\mathrm{Stage1}}
    = \mathcal{L}_{\mathrm{simple}}
    + \lambda_{\mathrm{div}}\, \mathcal{L}_{\mathrm{div}}
    + \lambda_{\mathrm{align}}\, \mathcal{L}_{\mathrm{align}},
    \quad \lambda_{\mathrm{div}}, \lambda_{\mathrm{align}} > 0.
\end{equation}
\end{proposition}

\begin{proof}
Suppose for contradiction that some global minimiser has
$\hat{s}_i(x_0) = c$ for all $i$ and some image $x_0$.

\textbf{Step 1 (Diversity).}
By Lemma~\ref{lem:div-collapse}(b), the collapsed configuration has
$\mathcal{L}_{\mathrm{div}}(x_0) = -d\log\epsilon$.
By assumption (A1), the teacher covariance $C^\star(x_0)$ has at least one
positive eigenvalue.
Setting $\hat{S}(x_0) = S^\star(x_0)$ yields
$C(x_0) = C^\star(x_0)$, which strictly decreases
$\mathcal{L}_{\mathrm{div}}(x_0)$ by Lemma~\ref{lem:div-collapse}(c).

\textbf{Step 2 (Reconstruction).}
Since $C^\star(x_0) \neq 0$, the teacher embeddings are not all equal,
so $\hat{S}(x_0) = (c, \ldots, c) \neq S^\star(x_0)$.
By (A2), this means $\mathcal{L}_{\mathrm{simple}}$ is not at its minimum,
and replacing $\hat{S}(x_0)$ by $S^\star(x_0)$ strictly reduces it.

\textbf{Step 3 (Alignment).}
By (A3) and Lemma~\ref{lem:align-collapse},
the collapsed configuration yields the maximal InfoNCE loss $\log B$,
which is strictly suboptimal.

\textbf{Step 4 (Contradiction).}
Since all three terms can be strictly decreased simultaneously
(by setting $\hat{S}(x_0) = S^\star(x_0)$), and $\lambda_{\mathrm{div}}, \lambda_{\mathrm{align}} > 0$,
the total loss $\mathcal{L}_{\mathrm{Stage1}}$ decreases,
contradicting the assumption of global optimality.
\end{proof}

\subsection{Partition-Based Spatial Modulation via SAM Masks}
\label{sec:partition}

We now formalise how the SAM partition induces spatially adaptive
conditioning, and show that cross-attention generalises this mechanism.

\begin{lemma}[Piecewise-constant modulation from a partition]
\label{lem:partition}
Let $\{M_i\}_{i=1}^{N}$ be a partition of the image grid $\Omega$
(i.e., $M_i \cap M_j = \emptyset$ for $i \neq j$ and
$\bigcup_i M_i = \Omega$), and let $s_i \in \mathbb{R}^{d}$ be
the region code for $M_i$.
Let $W_\gamma, W_\beta \in \mathbb{R}^{C \times d}$ be linear maps,
and define
$\gamma_i = W_\gamma s_i$, $\beta_i = W_\beta s_i \in \mathbb{R}^C$.
Set
\begin{equation}\label{eq:pw-modulation}
    \gamma_c(u,v) = \gamma_{i(u,v),c}, \quad
    \beta_c(u,v) = \beta_{i(u,v),c},
    \quad (u,v) \in \Omega,\; c = 1,\ldots,C,
\end{equation}
where $i(u,v)$ is the unique index with $(u,v) \in M_{i(u,v)}$.
Then:
\begin{enumerate}
    \item[\textup{(a)}] $\gamma$ and $\beta$ are piecewise constant on
    $\{M_i\}$; and
    \item[\textup{(b)}] conversely, any piecewise-constant spatial modulation
    field on $\{M_i\}$ that depends linearly on the region codes
    $\{s_i\}$ can be written in this form.
\end{enumerate}
\end{lemma}

\begin{proof}
(a) follows directly from the construction: for $(u,v), (u',v') \in M_i$,
we have $i(u,v) = i(u',v') = i$, so $\gamma_c(u,v) = \gamma_{i,c} = \gamma_c(u',v')$.

(b) Any piecewise-constant field assigns a vector $\tilde{\gamma}_i \in \mathbb{R}^C$
to each region $M_i$.
If this depends linearly on $s_i$, then
$\tilde{\gamma}_i = F(s_i) = W_\gamma s_i$ for some $W_\gamma \in \mathbb{R}^{C \times d}$.
\end{proof}

\begin{proposition}[Cross-attention generalises partition modulation]
\label{prop:cross-attn-generalises}
The cross-attention mechanism in Definition~\ref{def:seg-cross-attn}
is strictly more expressive than the piecewise-constant modulation
of Lemma~\ref{lem:partition}.
Specifically, by concentrating attention weights on a single token per
spatial location, cross-attention can recover the hard partition assignment
$(u,v) \mapsto i(u,v)$ as a special case, while also permitting
smooth blending across region boundaries.
\end{proposition}

\begin{proof}
In cross-attention, the output at spatial position $(u,v)$ is
\begin{equation}\label{eq:xattn-output}
    o_{(u,v)} = \sum_{i=1}^{N} a_{(u,v),i}\, V_i,
    \quad
    a_{(u,v),i} = \frac{\exp(Q_{(u,v)}^\top K_i / \sqrt{d/n_h})}
    {\sum_{j=1}^{N} \exp(Q_{(u,v)}^\top K_j / \sqrt{d/n_h})},
\end{equation}
where $Q_{(u,v)} = W_Q\, x_{(u,v)}$ and $K_i = V_i = W_S\, s_i$.

In the limit where attention weights concentrate on a single token
(i.e., $a_{(u,v),i} \to \mathbf{1}[i = i(u,v)]$), the output becomes
$o_{(u,v)} = W_S\, s_{i(u,v)}$, which is exactly the piecewise-constant
modulation with $W_\gamma = W_O W_S$.
Since the attention weights can also take intermediate values,
cross-attention is strictly more expressive.
\end{proof}

\begin{remark}[Connection to SPADE and spatially adaptive normalisation]
\label{rem:spade}
The partition-based modulation of Lemma~\ref{lem:partition} is
closely related to SPADE (Spatially-Adaptive Denormalisation)~\cite{park2019spade},
which applies spatially varying affine transforms conditioned on a
semantic layout.
Our cross-attention mechanism can be seen as a \emph{learned, soft}
variant of SPADE where the spatial layout emerges from attention
patterns rather than being provided as a hard label map.
\end{remark}

\subsection{Representation Sufficiency}
\label{sec:sufficiency}

We justify using $z = (g, S)$ as the latent variable in our
factorisation~\eqref{eq:marginalisation}.

\begin{definition}[Sufficiency for region prediction]
\label{def:sufficiency}
Let $\mathcal{R}$ denote the space of region descriptors (e.g., indices
or positional encodings of masked patch subsets in the I-JEPA framework).
For each image $x$ and region descriptor $R \in \mathcal{R}$,
let $y_R(x)$ denote the \emph{target encoder} output
(the EMA teacher ViT representation of the patch subset specified by $R$,
as in~\cite{assran2023ijepa}).
We say that a representation $z$ is \emph{sufficient for the region
prediction family} $\{y_R\}_{R \in \mathcal{R}}$ if, for all $R$,
$\mathbb{E}[y_R(X) \mid X, R] = \mathbb{E}[y_R(X) \mid z, R]$
almost surely.
\end{definition}

\begin{proposition}[Sufficiency of $(g, S)$]
\label{prop:sufficiency}
Assume:
\begin{enumerate}
    \item[\textup{(A1)}] The I-JEPA context encoder $\phi_{\mathrm{jepa}}$
    and predictor $P$ are jointly trained to minimise
    $\mathbb{E}_{X, R}[\lVert P(g, R) - y_R(X) \rVert^2]$
    with fixed target encoder outputs $y_R(X)$.
    \item[\textup{(A2)}] The predictor class is rich enough to attain the
    Bayes-optimal regression for a.e.\ $(g, R)$.
\end{enumerate}
Then: (i) the optimal predictor satisfies
$P^\star(g, R) = \mathbb{E}[y_R(X) \mid g, R]$; and
(ii) $\mathbb{E}[y_R(X) \mid g, R] = \mathbb{E}[y_R(X) \mid z, R]$,
so $z = (g, S)$ is sufficient in the sense of
Definition~\ref{def:sufficiency}.
\end{proposition}

\begin{proof}
(i) Under squared-error loss, the Bayes-optimal predictor is the
conditional expectation~\cite{lehmann2006theory}:
$P^\star(g, R) = \arg\min_{f} \mathbb{E}[\lVert f(g, R) - y_R(X) \rVert^2 \mid g, R]
= \mathbb{E}[y_R(X) \mid g, R]$.

(ii) Since $g = \phi_{\mathrm{jepa}}(X)$ and $S = \phi_{\mathrm{sam}}(X)$
are deterministic functions of $X$, the $\sigma$-algebra satisfies
$\sigma(g, R) \subseteq \sigma(z, R)$.
By the tower property,
$\mathbb{E}[y_R(X) \mid z, R]
= \mathbb{E}[\mathbb{E}[y_R(X) \mid g, R] \mid z, R]
= \mathbb{E}[y_R(X) \mid g, R]$,
since $\mathbb{E}[y_R(X) \mid g, R]$ is $\sigma(g, R)$-measurable
and hence $\sigma(z, R)$-measurable.
\end{proof}

\begin{remark}[Interpretation]
Proposition~\ref{prop:sufficiency} shows that $g$ alone carries
sufficient information for predicting I-JEPA's target encoder outputs.
The SAM component $S$ is therefore \emph{complementary}: it encodes
spatial-structural information (object boundaries, part decomposition)
that is useful for pixel-level generation but not strictly needed for
the I-JEPA prediction task.
Together, $(g, S)$ captures both global semantics and local structure.
\end{remark}

%% ====================================================================
\section{Inference Algorithm}
Check \ref{alg:sampling}
\label{sec:inference}
%% ====================================================================

\begin{algorithm}[t]
\caption{Two-stage sampling procedure}
\label{alg:sampling}
\begin{algorithmic}[1]
\REQUIRE Trained Stage-1 diffusion model $f_\theta$, trained Stage-2
generator $G_\psi$, DDIM steps $S$, stochasticity $\eta$
\STATE $z_T \sim \mathcal{N}(0, I)$ \hfill\COMMENT{Sample noise in representation space}
\FOR{$i = S, S-1, \ldots, 1$}
    \STATE $\hat{z}_0 \leftarrow f_\theta(z_{t_i}, t_i)$ \hfill\COMMENT{$x_0$-prediction}
    \STATE $\hat{\epsilon} \leftarrow (z_{t_i} - \sqrt{\bar{\alpha}_{t_i}}\hat{z}_0) / \sqrt{1 - \bar{\alpha}_{t_i}}$ \hfill\COMMENT{Implied noise}
    \STATE $z_{t_{i-1}} \leftarrow \text{DDIM-update}(z_{t_i}, \hat{z}_0, \hat{\epsilon}, \eta)$ \hfill\COMMENT{Eq.~\eqref{eq:ddim-update}}
\ENDFOR
\STATE $\hat{g} \leftarrow (z_0)_0$, \quad $\hat{S} \leftarrow (z_0)_{1:N}$ \hfill\COMMENT{Split global and segment tokens}
\STATE $\hat{x} \leftarrow G_\psi(\hat{g}, \hat{S})$ \hfill\COMMENT{Decode image via StyleGAN2}
\RETURN $\hat{x}$
\end{algorithmic}
\end{algorithm}

%% ====================================================================
\section{Summary of Architectural Choices}
\label{sec:summary}
%% ====================================================================

\begin{table}[t]
\centering
\caption{Summary of key design choices.}
\label{tab:summary}
\begin{tabular}{@{}lll@{}}
\toprule
\textbf{Component} & \textbf{Choice} & \textbf{Details} \\
\midrule
\multicolumn{3}{@{}l}{\textit{Stage~1: Representation Diffusion}} \\
Diffusion framework & DDPM~\cite{ho2020ddpm} & Discrete-time, Gaussian \\
Parameterisation & $x_0$-prediction & Predicts clean tokens directly \\
Noise schedule & Linear & $\beta \in [10^{-4}, 0.02]$, $T = 1000$ \\
Backbone & UnifiedSegTransformer & 8 layers, $d = 768$, 12 heads \\
Timestep conditioning & AdaLN~\cite{peebles2023dit} & Scale--shift from sinusoidal embedding \\
Sampler & DDIM~\cite{song2021ddim} & 50 steps, $\eta = 0$ \\
\midrule
\multicolumn{3}{@{}l}{\textit{Stage~2: Conditional Generator}} \\
Base architecture & StyleGAN2-ADA~\cite{karras2020training} & Residual synthesis \\
I-JEPA injection & $\mathcal{W}$-space additive shift & Gated, first 6 layers \\
SAM injection & Cross-attention & Res.\ $16\times16$, $32\times32$ \\
Gate initialisation & Zero & Starts unconditional \\
Adversarial loss & Non-saturating logistic & With R1 penalty \\
Feature matching & Cosine (I-JEPA) & With linear warmup \\
\bottomrule
\end{tabular}
\end{table}

\bibliographystyle{plain}
\begin{thebibliography}{99}

\bibitem{assran2023ijepa}
M.~Assran, Q.~Duval, I.~Misra, P.~Bojanowski, P.~Vincent, M.~Rabbat,
Y.~LeCun, and N.~Ballas.
\newblock Self-supervised learning from images with a joint-embedding
predictive architecture.
\newblock In \emph{CVPR}, 2023.

\bibitem{bachlechner2021rezero}
T.~Bachlechner, B.~P.~Majumder, H.~H.~Mao, G.~W.~Cottrell, and J.~McAuley.
\newblock {ReZero} is all you need: Fast convergence at large depth.
\newblock In \emph{UAI}, 2021.

\bibitem{bardes2022vicreg}
A.~Bardes, J.~Ponce, and Y.~LeCun.
\newblock {VICReg}: Variance-invariance-covariance regularization for
self-supervised learning.
\newblock In \emph{ICLR}, 2022.

\bibitem{bordes2022rdm}
F.~Bordes, E.~Pimentel, K.~Srinivasan, and A.~Romero-Soriano.
\newblock High fidelity visualization of what your self-supervised
representation knows.
\newblock \emph{arXiv:2112.09164}, 2022.

\bibitem{goodfellow2014gan}
I.~Goodfellow, J.~Pouget-Abadie, M.~Mirza, B.~Xu, D.~Warde-Farley,
S.~Ozair, A.~Courville, and Y.~Bengio.
\newblock Generative adversarial nets.
\newblock In \emph{NeurIPS}, 2014.

\bibitem{goyal2017accurate}
P.~Goyal, P.~Doll{\'a}r, R.~Girshick, P.~Noordhuis, L.~Wesolowski,
A.~Kyrola, A.~Tulloch, Y.~Jia, and K.~He.
\newblock Accurate, large minibatch {SGD}: Training {ImageNet} in 1~hour.
\newblock \emph{arXiv:1706.02677}, 2017.

\bibitem{hendrycks2016gelu}
D.~Hendrycks and K.~Gimpel.
\newblock {Gaussian} error linear units ({GELUs}).
\newblock \emph{arXiv:1606.08415}, 2016.

\bibitem{ho2020ddpm}
J.~Ho, A.~Jain, and P.~Abbeel.
\newblock Denoising diffusion probabilistic models.
\newblock In \emph{NeurIPS}, 2020.

\bibitem{ho2022cfg}
J.~Ho and T.~Salimans.
\newblock Classifier-free diffusion guidance.
\newblock In \emph{NeurIPS Workshop}, 2022.

\bibitem{karras2019style}
T.~Karras, S.~Laine, and T.~Aila.
\newblock A style-based generator architecture for generative adversarial
networks.
\newblock In \emph{CVPR}, 2019.

\bibitem{karras2020stylegan2}
T.~Karras, S.~Laine, M.~Aittala, J.~Hellsten, J.~Lehtinen, and T.~Aila.
\newblock Analyzing and improving the image quality of {StyleGAN}.
\newblock In \emph{CVPR}, 2020.

\bibitem{karras2020training}
T.~Karras, M.~Aittala, J.~Hellsten, S.~Laine, J.~Lehtinen, and T.~Aila.
\newblock Training generative adversarial networks with limited data.
\newblock In \emph{NeurIPS}, 2020.

\bibitem{kirillov2023sam}
A.~Kirillov, E.~Mintun, N.~Ravi, H.~Mao, C.~Rolland, L.~Gustafson,
T.~Xiao, S.~Whitehead, A.~C.~Berg, W.-Y.~Lo, P.~Doll{\'a}r, and
R.~Girshick.
\newblock Segment anything.
\newblock In \emph{ICCV}, 2023.

\bibitem{kulesza2012dpp}
A.~Kulesza and B.~Taskar.
\newblock Determinantal point processes for machine learning.
\newblock \emph{Foundations and Trends in Machine Learning}, 5(2--3):123--286,
2012.

\bibitem{lehmann2006theory}
E.~L.~Lehmann and G.~Casella.
\newblock \emph{Theory of Point Estimation}.
\newblock Springer, 2nd edition, 2006.

\bibitem{mescheder2018training}
L.~Mescheder, A.~Geiger, and S.~Nowozin.
\newblock Which training methods for {GANs} do actually converge?
\newblock In \emph{ICML}, 2018.

\bibitem{miyato2018projection}
T.~Miyato and M.~Koyama.
\newblock c{GAN}s with projection discriminator.
\newblock In \emph{ICLR}, 2018.

\bibitem{oord2018infonce}
A.~van~den~Oord, Y.~Li, and O.~Vinyals.
\newblock Representation learning with contrastive predictive coding.
\newblock \emph{arXiv:1807.03748}, 2018.

\bibitem{park2019spade}
T.~Park, M.-Y.~Liu, T.-C.~Wang, and J.-Y.~Zhu.
\newblock Semantic image synthesis with spatially-adaptive normalization.
\newblock In \emph{CVPR}, 2019.

\bibitem{peebles2023dit}
W.~Peebles and S.~Xie.
\newblock Scalable diffusion models with transformers.
\newblock In \emph{ICCV}, 2023.

\bibitem{ramesh2022hierarchical}
A.~Ramesh, P.~Dhariwal, A.~Nichol, C.~Chu, and M.~Chen.
\newblock Hierarchical text-conditional image generation with {CLIP} latents.
\newblock \emph{arXiv:2204.06125}, 2022.

\bibitem{sohl2015deep}
J.~Sohl-Dickstein, E.~A.~Weiss, N.~Maheswaranathan, and S.~Ganguli.
\newblock Deep unsupervised learning using nonequilibrium thermodynamics.
\newblock In \emph{ICML}, 2015.

\bibitem{song2021ddim}
J.~Song, C.~Meng, and S.~Ermon.
\newblock Denoising diffusion implicit models.
\newblock In \emph{ICLR}, 2021.

\bibitem{vaswani2017attention}
A.~Vaswani, N.~Shazeer, N.~Parmar, J.~Uszkoreit, L.~Jones, A.~N.~Gomez,
\L.~Kaiser, and I.~Polosukhin.
\newblock Attention is all you need.
\newblock In \emph{NeurIPS}, 2017.

\end{thebibliography}

\end{document}
